\documentclass[12pt,a4paper]{article}
\usepackage[utf8]{inputenc}
\usepackage{geometry}
\geometry{a4paper, margin=1in}
\usepackage{graphicx}
\usepackage{hyperref}
\usepackage{float}
\usepackage{setspace}

\title{\textbf{Software Design - Assignment 1\\Customer Relationship Management (CRM)}}
\author{Computer Science Student}
\date{\today}

\begin{document}

\maketitle
\newpage

\tableofcontents
\newpage

\section{Introduction}
\spacing{1.5}
This is the documentation for the first assignment in Software Design. For this project I choosed to make a Customer Relationship Management (CRM) system because it seemed like a practical idea. Basically, the goal was to manage a domain entity, so I picked the Customer entity.

Any shop or business needs a CRM to keep track of the clients and all the \textbf{informations} related to them. This project is a simple web application that let us do just that, using basic object-oriented principles.

\section{Problem Definition}
The app is pretty straightforward: it manages Customers. A Customer is basically someone who buys things from us, and we need to save their name, email, phone, and company.

Also, the assignment PDF said we need roles, so I added three of them:
\begin{itemize}
    \item \textbf{Visitor} - someone who just opens the site. They can only see the list, nothing else.
    \item \textbf{Administrator} - the boss. Can do everything, including deleting customers.
    \item \textbf{Employee} - they need to login. They can add new customers or edit the mistakes, but they can't delete anything by mistake.
\end{itemize}

\section{Application Architecture}
I tried to follow the MVC (Model-View-Controller) architecture exactly like it was required. The code is written in Java and I used Spring Boot because it does a lot of the heavy lifting.

\subsection{Model}
Here we have the entities and the repositories. For the database I didn't want to complicate myself with installing MySQL on my laptop and configuring users, so I used an H2 in-memory database. It works out of the box with Spring Data JPA.

\subsection{View}
For the frontend I just used standard HTML pages with Thymeleaf. I also threw in some Bootstrap classes so it doesn't look like it was made in the 90s. It has a start page, a login, and the tables.

\subsection{Controller}
The controllers take the requests from the browser and tell the app what to do. The controller \textbf{is depending of} the service layer, so there is no direct database code in the controller. I think this is what "clean architecture" means.

\section{Functional Requirements}
\subsection{CRUD Operations}
You can do all the basic CRUD stuff. 
\begin{itemize}
    \item \textbf{Create}: If you are logged in as employee or admin, you have a form to add a new customer (Name, Email, Phone, Company).
    \item \textbf{Read}: Visitors can see the table with the clients.
    \item \textbf{Update}: You can click "Edit" to fix a typo in a customer's email.
    \item \textbf{Delete}: The delete button only shows up if you are logged in as admin.
\end{itemize}

\subsection{Extra Features}
I also added a search bar on top of the list. If you type a name, it filters the table. I thought this was the easiest extra feature to implement to get those 2 points.

\section{Security}
I used Spring Security to handle the logins. The passwords are encrypted with BCrypt so they are safe in the database. If you are not logged in, spring security just treats you as a Visitor.

\section{Diagrams}
Here are the diagrams requested in the deliverables section. 

\subsection{Use Case Diagram}
This shows who can do what. Admin can do everything, Employee can do almost everything, Visitor just watches.

\begin{figure}[H]
    \centering
    \includegraphics[width=0.8\textwidth]{diagrams/usecase.png}
    \caption{Use Case Diagram}
\end{figure}

\subsection{Entity-Relationship Diagram}
I only have one table right now, called \texttt{customers}, with an auto-increment ID.

\begin{figure}[H]
    \centering
    \includegraphics[width=0.4\textwidth]{diagrams/er.png}
    \caption{Entity-Relationship Diagram}
\end{figure}

\subsection{Class Diagram}
Here we can see the classes I wrote in Java. \texttt{CustomerController} talks to \texttt{CustomerService}, which talks to \texttt{CustomerRepository}.

\begin{figure}[H]
    \centering
    \includegraphics[width=0.9\textwidth]{diagrams/class.png}
    \caption{Class Diagram}
\end{figure}

\subsection{Sequence Diagram}
This is the flow when someone clicks "Save Customer" in the UI form. It travels through all the layers and saves in H2.

\begin{figure}[H]
    \centering
    \includegraphics[width=\textwidth]{diagrams/sequence.png}
    \caption{Sequence Diagram: Add New Customer}
\end{figure}

\section{Conclusions}
To sum up, the application works and all the required features from the PDF are done. Using Spring Boot was a good idea because it saved me a lot of time on the configuration part. The project is ready to be expanded for Assignment 2.

\end{document}
